\centerline{\bf \Huge Домашнее задание}
\centerline{\bf \Huge Лемма о венике}

\begin{flushright}
        \scriptsize{Нет никакого смысла трудиться в поте лица,\\ если ты не веришь в себя.}
\end{flushright}

\problem{1}
Вокруг прямоугольного треугольника $A B C$ с прямым углом $C$ описана окружность, на меньших дугах $A C$ и $B C$ взяты их середины --- $K$ и $P$ соответственно. Отрезок $K P$ пересекает катет $A C$ в точке $N$. Центр вписанной окружности треугольника $A B C$ --- $I$. Найти угол $NIC$.

\problem{2}Точки $A, B$ и $C$ лежат на окружности с центром в точке $O$. Луч $O B$ вторично пересекает описанную около треугольника $A O C$ в точке $D$, причём точка $B$ оказалась внутри этой окружности. Докажите, что $A B$ --- биссектриса угла $D A C$.

\problem{3}Биссектрисы углов $B$ и $C$ треугольника $A B C$ пересекают его описанную окружность в точках $B_0$ и $C_0$ соответственно и пересекаются в точке $I$. Окружности $\omega_B$ и $\omega_C$ с центрами $B_0$ и $C_0$ касаются отрезков $A C$ и $A B$ соответственно. Докажите, что через точку $I$ можно провести прямую, касающуюся обеих окружностей $\omega_B$ и $\omega_C$.

\problem{4}Дан треугольник $A B C$. На продолжениях сторон $A B$ и $C B$ за точку $B$ взяты точки $C_1$ и $A_1$ соответственно так, что $A C=A_1 C=A C_1$. Докажите, что описанные окружности треугольников $A B A_1$ и $C B C_1$ пересекаются на биссектрисе угла $B$.

\problem{5}
Биссектрисы углов $A$ и $C$ треугольника $A B C$ пересекают описанную окружность этого треугольника в точках $A_0$ и $C_0$ соответственно. Прямая, проходящая через центр вписанной окружности треугольника $A B C$ параллельно стороне $A C$, пересекается с прямой $A_0 C_0$ в точке $P$. Докажите, что прямая $P B$ касается описанной окружности треугольника $A B C$.

\problem{6}
Точка $I$ - центр вписанной окружности треугольника $A B C$, точка $M$ - середина стороны $A C$, а точка $W$ - середина не содержащей $C$ дуги $A B$ описанной окружности. Оказалось, что $\angle A I M=90^{\circ}$. В каком отношении $I$ делит отрезок $C W$ ?

\problem{7}
На дугах $A B$ и $B C$ окружности, описанной около треугольника $A B C$, выбраны соответственно точки $K$ и $L$ так, что прямые $K L$ и $A C$ параллельны. Докажите, что центры вписанных окружностей треугольников $A B K$ и $C B L$ равноудалены от середины дуги $A B C$.